% GENERATED FILE. Edit canonical lessons, then run npm run generate:books.
% canonical-source-hash: fnv1a64:36f0cb8ef1e20044
% canonical-lessons: TA-C144-uyaramana, TA-C144-cutana, TA-C144-kulirnta, TA-C144-kanamana, TA-C144-lecana

\chapter{Tall, Hot, Cold, Heavy, Light}
\label{ch:ta-tall-hot-cold-heavy-light}
\hlchaptermodality{144}

\begin{chapteropening}
\textbf{By the end of this chapter:} \emph{I can say tall, hot, cold, heavy and light.}

Complete the last lesson of chapter 144: I can say tall, hot, cold, heavy and light.
\end{chapteropening}

\section[\texorpdfstring{\ta{உயரமான} (uyaramāṉa)}{உயரமான (uyaramāṉa)}]{\ta{உயரமான} (uyaramāṉa) — tall}

\label{lesson:TA-C144-uyaramana}



\begin{quote}
Before the new one: say the Tamil for long, then the Tamil for short.
\end{quote}

\subsection*{You'll want to know: \ta{உயரமான}}
\textbf{\ta{உயரமான}} — \emph{uyaramāṉa} — \textquotedblleft{}tall\textquotedblright{}.

\textbf{In the family, and next door}

\noindent
\begin{tabularx}{\linewidth}{@{}>{\raggedright\arraybackslash}X>{\raggedright\arraybackslash}X@{}}
\toprule
\textbf{} & \textbf{word} \\
\midrule
Telugu & \te{పొడవు} (\emph{poḍavu}) \\
Hindi (neighbour) & \dv{लंबा} (\emph{lambā}) \\
English & tall \\
\bottomrule
\end{tabularx}

\begin{grammarlens}[title={What you've built}]
One more word for this chapter.
\end{grammarlens}

\subsection*{Your turn}
\begin{itemize}
\raggedright
  \item \emph{Say it:} \emph{uyaramāṉa}
  \item \emph{Say it:} \emph{uyaramāṉa}, once more
  \item \emph{Say it:} say \emph{uyaramāṉa}
  \item \emph{Recall:} say the Tamil for long, then the Tamil for short, then say \emph{uyaramāṉa} again
\end{itemize}

\begin{tcolorbox}[breakable,colback=teal!4,colframe=teal!35!black,title={Before you move on}]
What does \ta{உயரமான} mean? (\textquotedblleft{}Tall\textquotedblright{}.) Say it once more.
\end{tcolorbox}
\section[\texorpdfstring{\ta{சூடான} (cūṭāṉa)}{சூடான (cūṭāṉa)}]{\ta{சூடான} (cūṭāṉa) — hot}

\label{lesson:TA-C144-cutana}



\begin{quote}
Before the new one: say the Tamil for short, then the Tamil for tall.

Tell the time, \textbf{\ta{ஒரு} \ta{மணி}}, \textbf{\ta{இரண்டு} \ta{மணி}}, then name the afternoon three ways: \textbf{\ta{நடுப்பகல்}}, \textbf{\ta{மதியம்}}, \textbf{\ta{பிற்பகல்}}.
\end{quote}

\subsection*{You'll want to know: \ta{சூடான}}
\textbf{\ta{சூடான}} — \emph{cūṭāṉa} — \textquotedblleft{}hot\textquotedblright{}.

\textbf{In the family, and next door}

\noindent
\begin{tabularx}{\linewidth}{@{}>{\raggedright\arraybackslash}X>{\raggedright\arraybackslash}X>{\raggedright\arraybackslash}X@{}}
\toprule
\textbf{} & \textbf{word} & \textbf{same root} \\
\midrule
Malayalam & \ml{ചൂടുള്ള} (\emph{cūṭuḷḷa}) & yes \\
Hindi (neighbour) & \dv{गरम} (\emph{garam}) &  \\
English & hot &  \\
\bottomrule
\end{tabularx}

\begin{grammarlens}[title={What you've built}]
One more word for this chapter.
\end{grammarlens}

\subsection*{Your turn}
\begin{itemize}
\raggedright
  \item \emph{Say it:} \emph{cūṭāṉa}
  \item \emph{Say it:} \emph{cūṭāṉa}, once more
  \item \emph{Say it:} say \emph{cūṭāṉa}
  \item \emph{Recall:} say the Tamil for short, then the Tamil for tall, then say \emph{cūṭāṉa} again
\end{itemize}

\begin{tcolorbox}[breakable,colback=teal!4,colframe=teal!35!black,title={Before you move on}]
What does \ta{சூடான} mean? (\textquotedblleft{}Hot\textquotedblright{}.) Say it once more.
\end{tcolorbox}
\section[\texorpdfstring{\ta{குளிர்ந்த} (kuḷirnta)}{குளிர்ந்த (kuḷirnta)}]{\ta{குளிர்ந்த} (kuḷirnta) — cold}

\label{lesson:TA-C144-kulirnta}



\begin{quote}
Before the new one: say the Tamil for tall, then the Tamil for hot.
\end{quote}

\subsection*{You'll want to know: \ta{குளிர்ந்த}}
\textbf{\ta{குளிர்ந்த}} — \emph{kuḷirnta} — \textquotedblleft{}cold\textquotedblright{}.

\textbf{In the family, and next door}

\noindent
\begin{tabularx}{\linewidth}{@{}>{\raggedright\arraybackslash}X>{\raggedright\arraybackslash}X@{}}
\toprule
\textbf{} & \textbf{word} \\
\midrule
Kannada & \ka{ಶೀತ} (\emph{śīta}) \\
English & cold \\
\bottomrule
\end{tabularx}

\begin{grammarlens}[title={What you've built}]
One more word for this chapter.
\end{grammarlens}

\subsection*{Your turn}
\begin{itemize}
\raggedright
  \item \emph{Say it:} \emph{kuḷirnta}
  \item \emph{Say it:} \emph{kuḷirnta}, once more
  \item \emph{Say it:} say \emph{kuḷirnta}
  \item \emph{Recall:} say the Tamil for tall, then the Tamil for hot, then say \emph{kuḷirnta} again
\end{itemize}

\begin{tcolorbox}[breakable,colback=teal!4,colframe=teal!35!black,title={Before you move on}]
What does \ta{குளிர்ந்த} mean? (\textquotedblleft{}Cold\textquotedblright{}.) Say it once more.
\end{tcolorbox}
\section[\texorpdfstring{\ta{கனமான} (kaṉamāṉa)}{கனமான (kaṉamāṉa)}]{\ta{கனமான} (kaṉamāṉa) — heavy}

\label{lesson:TA-C144-kanamana}



\begin{quote}
Before the new one: say the Tamil for hot, then the Tamil for cold.
\end{quote}

\subsection*{You'll want to know: \ta{கனமான}}
\textbf{\ta{கனமான}} — \emph{kaṉamāṉa} — \textquotedblleft{}heavy\textquotedblright{}.

\textbf{In the family, and next door}

\noindent
\begin{tabularx}{\linewidth}{@{}>{\raggedright\arraybackslash}X>{\raggedright\arraybackslash}X@{}}
\toprule
\textbf{} & \textbf{word} \\
\midrule
Kannada & \ka{ಭಾರ} (\emph{bhāra}) \\
Telugu & \te{బరువు} (\emph{baruvu}) \\
Malayalam & \ml{ഭാരമുള്ള} (\emph{bhāramuḷḷa}) \\
Hindi (neighbour) & \dv{भारी} (\emph{bhārī}) \\
English & heavy \\
\bottomrule
\end{tabularx}

\begin{grammarlens}[title={What you've built}]
One more word for this chapter.
\end{grammarlens}

\subsection*{Your turn}
\begin{itemize}
\raggedright
  \item \emph{Say it:} \emph{kaṉamāṉa}
  \item \emph{Say it:} \emph{kaṉamāṉa}, once more
  \item \emph{Say it:} say \emph{kaṉamāṉa}
  \item \emph{Recall:} say the Tamil for hot, then the Tamil for cold, then say \emph{kaṉamāṉa} again
\end{itemize}

\begin{tcolorbox}[breakable,colback=teal!4,colframe=teal!35!black,title={Before you move on}]
What does \ta{கனமான} mean? (\textquotedblleft{}Heavy\textquotedblright{}.) Say it once more.
\end{tcolorbox}
\section[\texorpdfstring{\ta{லேசான} (lēcāṉa)}{லேசான (lēcāṉa)}]{\ta{லேசான} (lēcāṉa) — light (in weight)}

\label{lesson:TA-C144-lecana}



\begin{quote}
Before the new one: say the Tamil for cold, then the Tamil for heavy.
\end{quote}

\subsection*{You'll want to know: \ta{லேசான}}
\textbf{\ta{லேசான}} — \emph{lēcāṉa} — \textquotedblleft{}light (in weight)\textquotedblright{}.

\textbf{In the family, and next door}

\noindent
\begin{tabularx}{\linewidth}{@{}>{\raggedright\arraybackslash}X>{\raggedright\arraybackslash}X@{}}
\toprule
\textbf{} & \textbf{word} \\
\midrule
Kannada & \ka{ಹಗುರ} (\emph{hagura}) \\
Telugu & \te{తేలిక} (\emph{tēlika}) \\
Hindi (neighbour) & \dv{हल्का} (\emph{halkā}) \\
English & light \\
\bottomrule
\end{tabularx}

\begin{grammarlens}[title={What you've built}]
One more word for this chapter.
\end{grammarlens}

\subsection*{Your turn}
\begin{itemize}
\raggedright
  \item \emph{Say it:} \emph{lēcāṉa}
  \item \emph{Say it:} \emph{lēcāṉa}, once more
  \item \emph{Say it:} say \emph{lēcāṉa}
  \item \emph{Recall:} say the Tamil for cold, then the Tamil for heavy, then say \emph{lēcāṉa} again
\end{itemize}

\begin{tcolorbox}[breakable,colback=teal!4,colframe=teal!35!black,title={Before you move on}]
What does \ta{லேசான} mean? (\textquotedblleft{}Light (in weight)\textquotedblright{}.) Say it once more.
\end{tcolorbox}
